Las mediciones experimentales se realizaron en un entorno Linux (Ubuntu 24.04.1 LTS sobre WSL2, arquitectura x86\_64), en una máquina equipada con un procesador Intel Core i7-1255U (10 núcleos, 12 hilos, con frecuencias de hasta 4.70~GHz) y 6.6~GiB de memoria RAM asignados al entorno virtualizado. Las soluciones fueron desarrolladas en C++17 y compiladas mediante \texttt{g++} versión 13.3.0 aplicando la bandera de optimización \texttt{-O3}. La medición temporal se instrumentó con \texttt{std::chrono::steady\_clock} con resolución de microsegundos ($\mu\text{s}$), mientras que el consumo de memoria máxima residente se obtuvo a través de la llamada al sistema \texttt{getrusage} (\texttt{ru\_maxrss}). Para evitar contaminación acumulativa de memoria entre corridas, cada ejecución individual se ejecutó como un proceso independiente orquestado desde el \texttt{makefile}.

Los experimentos se evaluaron sobre las instancias de prueba parametrizadas en el enunciado:
\begin{itemize}
    \item \textbf{Ordenamiento de Arreglos:} Arreglos de dimensión $n \in \{10^1, 10^3, 10^5, 10^7\}$, bajo tres tipos de orden previo (\emph{aleatorio}, \emph{ascendente} y \emph{descendente}) y dos rangos de dominio: $D_1 = \{0,\dots,9\}$ y $D_7 = \{0,\dots,10^7\}$, evaluando 3 muestras independientes ($a, b, c$) por cada combinación (288 ejecuciones por algoritmo).
    \item \textbf{Multiplicación de Matrices:} Pares de matrices cuadradas de orden $n \in \{2^4, 2^6, 2^8, 2^{10}\}$, abarcando configuraciones \emph{densas}, \emph{diagonales} y \emph{dispersas} (10\% de elementos no nulos), con coeficientes pertenecientes a los dominios $D_0 = \{0,1\}$ y $D_{10} = \{0,\dots,9\}$, considerando 3 repeticiones independientes por instancia.
\end{itemize}